\begin{abstract}
\noindent
Running a 27B-parameter model at home means choosing a quantization, a KV-cache precision, a context
length, a GPU split and a speculative decoder, usually from tables measured on other hardware
under an unstated protocol. This report measures those choices for Qwen3.8-27B on two 16\,GB
RTX~5060~Ti cards with 14\,GiB of system RAM, and reports what each instrument could and could not see.
\emph{Accuracy.} Benchmarks bound quantization damage but do not rank it: HellaSwag ($n{=}400$),
HumanEval+ ($n{=}164$), SWE-bench Verified ($n{\approx}50$), GPQA Diamond ($n{=}198$) and AA-LCR
($n{=}100$) all fail to separate a 4-bit from a 6-bit build, while 2.3 GPU-hours of KL divergence
separates every adjacent pair at 3.7--11.8\,$\sigma$. The damage is about twice as large
on code as on prose and sits in a thin tail of tokens. Speculative decoding is not output-identical
to plain decoding here: at temperature 0 with a fixed seed it changes 33 of 164 completions,
reproducibly, while task accuracy does not move.
\emph{Context.} The native 262{,}144-token window fits on 32\,GB of consumer VRAM only after
rebalancing the layer split by hand; the engine default fails for three of four builds, and the best
split does not transfer between builds. The 4-bit build holds that window with an 8-bit cache, the
6-bit build only with a 4-bit cache; both drafters reach the same window unless the cache
is unquantized.
\emph{Speed.} At the full window a smaller build is not measurably faster. On one engine build, a
block-diffusion drafter (DFlash2) at its design depth decodes 45--65\,\% faster than the built-in
multi-token-prediction head at every depth to 246K tokens, about $7\times$ faster than no speculation,
with 3.6\,GiB less host RAM. Host RAM, not VRAM, caused the only serving outage.
Sixteen of the study's own claims were withdrawn or narrowed on re-analysis, fourteen at no GPU cost
because per-item records were kept. All data, harness code and the decision log are public.
\end{abstract}
